
\subsubsection{Multi-Dimensional Evaluation Frameworks}

HELM (Holistic Evaluation of Language Models) aggregates performance across 50+ benchmarks spanning accuracy, robustness, fairness, bias, toxicity, and efficiency. While HELM represents comprehensive multi-dimensional evaluation, it reports separate scores per dimension without testing whether failures correlate or cluster into shared failure modes. The present work extends aggregation-based approaches by analyzing cross-dimensional correlation structure. Our finding that three HELM-included benchmarks (TrustfulQA, AdvBench, BOLD) correlate at r > 0.99 suggests score aggregation may mask latent structure. Where HELM treats benchmarks as independent scorecards, we treat them as measurement instruments for discovering correlation patterns.

BIG-bench aggregates 200+ tasks but focuses on capability breadth rather than trustworthiness-specific dimensions. Recent analyses identify task clustering based on difficulty but not cross-dimensional trustworthiness correlations. Our approach differs by explicitly testing the independence assumption: if reliability, robustness, and fairness arise from independent mechanisms, correlations should be near-zero. Observed r > 0.99 challenges this framing.

\subsubsection{Single-Dimension Trustworthiness Benchmarks}

\textbf{TrustfulQA} measures reliability by testing whether models generate truthful responses to 817 questions designed to elicit common misconceptions. The benchmark establishes that larger models do not necessarily produce more truthful outputs. Our work complements this by showing TrustfulQA scores correlate r = 0.998 with AdvBench and r = 0.993 with BOLD, suggesting reliability failures co-occur with robustness and fairness failures at near-identical rates.

\textbf{AdvBench} evaluates robustness to adversarial attacks through jailbreak prompts and input perturbations. Prior work documents that adversarially robust models often require explicit robustness training. Our correlation analysis shows AdvBench scores tightly coupled with TrustfulQA (r = 0.998) and BOLD (r = 0.997), raising the question whether robustness training simultaneously affects reliability and fairness---a prediction testable through targeted intervention experiments deferred to future work.

\textbf{BOLD} (Bias in Open-Ended Language Generation) quantifies fairness by measuring sentiment and regard differences across demographic groups. BOLD establishes that pre-trained models exhibit systematic bias amplification persisting after alignment training. Our finding of r > 0.99 correlations between BOLD and reliability/robustness benchmarks suggests bias amplification may share root causes with hallucination and adversarial brittleness, though the 3-benchmark design cannot distinguish unified constructs from insufficient measurement resolution.

\subsubsection{Trustworthiness Interventions}

Prior work documents that alignment fine-tuning improves multiple dimensions simultaneously. RLHF reduces hallucinations (reliability) while decreasing toxic outputs (safety) and improving demographic parity (fairness). However, this literature lacks quantitative taxonomy predicting which dimensions co-vary under interventions. Our correlation-based approach formalizes this intuition: if TrustfulQA and AdvBench correlate at r = 0.998, interventions targeting reliability plausibly affect robustness as well. Testing this prediction requires validated failure mode clusters, which our clustering analysis failed to produce (silhouette = 0.274 < 0.5), leaving intervention validation as future work.

\subsubsection{Positioning}

We extend multi-dimensional evaluation from aggregation (HELM) to correlation analysis, revealing that three major benchmarks measure near-redundant constructs (r > 0.99) rather than orthogonal dimensions. This challenges the design assumption in single-dimension benchmarks (TrustfulQA, AdvBench, BOLD), which were developed in isolation without testing empirical independence. Unlike intervention studies that qualitatively observe multi-dimensional improvements, we formalize correlation structure to enable mechanism-guided intervention design.

The key difference from prior work: we explicitly test whether trustworthiness dimensions are independent (null hypothesis: r $\approx$ 0) or correlated (alternative: r > 0.3), finding correlations $3\times$ stronger than hypothesized thresholds. This shifts evaluation paradigm from dimension-independent scorecards to correlation-aware diagnostic tools.
